% Transcription — fonds Alexandre Grothendieck, shelfmark 138, batch 4 (pages 61-80).
% Established from the facsimile archives/batches/138/batch-04.pdf (cover sheet
% excluded; batch page 1 = archivists' page 61), per the skill
% transcribe-grothendieck.
%
% Pass: Opus 5.5 (claude-opus-5-5), 2026-09-30 — first pass, unchecked against the pages by a human.
%
% Skipped: pp. 66 and 68 (typed USTL notices of thesis defences, June 1978,
% nothing in his hand). Pp. 64 and 69 are cover sheets carrying only a title in
% his hand; p. 75 likewise. P. 78 is blank. Cover words are the section
% headings below, with no \page.
\documentclass[11pt,a4paper]{article}
\input{../preamble/grothendieck}

\folder{138}
\batch{4}
\pages{61}{80}
\foldertitle{Cartes. Etude arithmétique (1976, 77 ?) : bon de commande (s.d.), notes manuscrites (s.d.).}
\dating{[vers 1976-1978]}
\watermark{Édition de démonstration}

\begin{document}

\page{61}
\note{la page s'ouvre sur la fin d'un calcul commencé avant le lot : $C$, $a$, $C_0$ et les $\pi_x$ n'y sont pas définis.}
On trouve
\[
  \operatorname{Aut}^{\circ}(C, a) \simeq \prod_{x \in C_0 \smallsetminus \{a\}} \pi_x
\]
\note{avant $\pi_x$, un symbole biffé, illisible.}

\drawing{un petit croquis : deux traits horizontaux coupés par quatre traits verticaux.}

\struck{Soit \ill{} $K$ \uncertain{des graphes}}

Soit $I$ un ensemble muni d'une structure polygonale d'ordre $n \geqslant 2$
\struck{i.e. \ill{}} [pour $n = 2$, il faut se donner l'un des
\add{de cardinal 2} \uncertain{orientés} $\Delta^{*}$, \struck{\ill{}}
\struck{\ill{} 2, et} \uncertain{ce qui revient au même},
$\Omega(\Pi) \simeq I \wedge I^{*}$ \ldots\ ]),
$D$ $(\simeq D_n)$ le groupe des automorphismes de $I$.

Soit $K = K(I)$ le graphe \uncertain{défini} par
\begin{itemize}
  \item \emph{sommets} : $\Omega(\Pi) = C_0 = \{o, o'\}$
  \item \emph{arêtes} : $I$ (indexation $(\ell_i)_{i \in I}$)
\end{itemize}
\struck{\ill{} est incidente}

Incidence : $C_0 \times I$ (\uncertain{tt} arête incidente à \uncertain{tt} sommet).

\drawing{le graphe $K(I)$ : un cercle, le sommet $o$ en haut, $o'$ en bas ; entre eux plusieurs arêtes parallèles marquées $\ell_0$, $\ell_1$, $\ell_2$.}

$G$ opère sur $K(I)$, et on a aussi \struck{\ill{}} \add{\ill{}} qui est
l'identité sur $I$, et \uncertain{échange} $o, o'$, commutant à $G$.
Soit $C(I)$ le groupe \uncertain{fondamental} \add{\ill{}} de $K(I)$.
$G$ y opère par automorphismes, en opérant transitivement sur
$C_0 = \{o, o'\}$.
On va regarder le \struck{\uncertain{commutant}} \add{normalisateur de $G$} dans

\page{62}
$G$ dans $\operatorname{Aut}(C(I))$. On a un élément canonique $\tau$
\note{la page s'arrête là ; le reste du feuillet est vide.}

\page{63}
\note{feuillet de papier teinté, sans texte : une figure et des formules.}
\drawing{une double pyramide : sommet $E^{+}$ en haut, $E^{-}$ en bas ;
à l'équateur le triangle $D_{20}$, $D_{01}$, $D_{12}$ (ce dernier en arrière).
Les arêtes sont les applications $p^{\pm}_{20}$, $p^{\pm}_{01}$, $p^{\pm}_{12}$
de $E^{\pm}$ vers les $D_{ij}$. Sur les faces sont placés $S_0$, $S_1$, $S_2$,
atteints par des flèches $\pi^{\pm}_0$, $\pi^{\pm}_1$, $\pi^{\pm}_2$ issues de
$E^{\pm}$, et par des flèches $q_{02}\colon D_{20} \to S_0$,
$q_{01}\colon D_{01} \to S_0$, $q_{20}$, $q_{10}$, $q_{12}$, $q_{21}$ issues
des sommets de l'équateur.}

\[
  \begin{cases}
    \rho_0^{+} = (p_{01}^{+})^{-1}\, p_{01}^{-}\, (p_{20}^{-})^{-1}\, p_{20}^{+} \\
    \rho_1^{+} = (p_{12}^{+})^{-1}\, p_{12}^{-}\, (p_{01}^{-})^{-1}\, p_{01}^{+} \\
    \rho_2^{+} = (p_{20}^{+})^{-1}\, p_{20}^{-}\, (p_{12}^{-})^{-1}\, p_{12}^{+}
  \end{cases}
\]
\[
  \rho_2^{+} \rho_1^{+} \rho_0^{+} = \mathrm{id}_{E^{+}}
\]
\note{dans la première ligne, les indices du premier facteur sont surchargés ; on les rétablit sur le modèle des deux lignes suivantes.}
\[
  \begin{aligned}
    S_0 &\simeq E^{+}/\rho_0^{+} \simeq E^{-}/\rho_0^{-} \\
    S_1 &\simeq E^{+}/\rho_1^{+} \simeq E^{-}/\rho_1^{-} \\
    S_2 &\simeq E^{+}/\rho_2^{+} \simeq E^{-}/\rho_2^{-}
  \end{aligned}
\]
\[
  \begin{cases}
    \rho_0^{-} = (p_{01}^{-})^{-1}\, p_{01}^{+}\, (p_{20}^{+})^{-1}\, p_{20}^{-} \\
    \rho_1^{-} = (p_{12}^{-})^{-1}\, p_{12}^{+}\, (p_{01}^{+})^{-1}\, p_{01}^{-} \\
    \rho_2^{-} = (p_{20}^{-})^{-1}\, p_{20}^{+}\, (p_{12}^{+})^{-1}\, p_{12}^{-}
  \end{cases}
\]

\section{Sous-groupes finis de $GP(1)$ (cas \emph{tame})}
\note{titre de sa main sur un feuillet par ailleurs blanc (p.~64), qui porte aussi, à gauche, un « S. » isolé.}

\page{65}
\uncertain{en général} \struck{$d_{ij} =$} $d_i = N/\nu_i$
\[
  2 - 2g' = N(2 - 2g) - \underbrace{(N\nu - \nu')}_{\sum_{i,j} (d_{ij} - 1)}
\]
\[
  \phantom{2 - 2g'} = N(2 - 2g) - \sum_i \underbrace{\nu_i (N/\nu_i - 1)}_{N - \nu_i \,=\, N(1 - \frac{\nu_i}{N})}
\]
\note{deux mots biffés sous le membre de gauche, illisibles.}
\[
  \boxed{\,2 - 2g' = N \Bigl[ (2 - 2g) - \sum_i \Bigl(1 - \frac{1}{d_i}\Bigr) \Bigr]\,}
\]
Si $g' = 0$, on trouve $g = 0$ et
\[
  \boxed{\,\frac{2}{N} = 2 - \sum_i \Bigl(1 - \frac{1}{d_i}\Bigr) > 0\,}
\]
d'où
\[
  \sum_i \Bigl(1 - \frac{1}{d_i}\Bigr) < 2 .
\]

\textbf{Obs.} Tous les $d_i \geqslant 2$ \struck{\ill{}} car $I$ \uncertain{contient}
\uncertain{l'ens. de ramification}.

1) $\operatorname{Card} I = 0$ : $\frac{2}{N} = 2$, $N = 1$ : \struck{pour} $G = 1$.

2) $\operatorname{Card} I = 1$ :
\[
  \frac{2}{N} = 2 - \Bigl(1 - \frac{1}{d}\Bigr) = 1 + \frac{1}{d}, \qquad
  2 = N + \frac{N}{d} = N + \nu
\]
\[
  N \geqslant 1,\ \nu \geqslant 1 \quad\text{donc}\quad N = 1,\ \nu = 1,\ I = \varnothing
\]
\emph{impossible}.

3) $\operatorname{Card} I = 2$ :
\[
  \frac{2}{N} = \frac{1}{d_1} + \frac{1}{d_2}, \qquad
  2 = \frac{N}{d_1} + \frac{N}{d_2} = \nu_1 + \nu_2, \qquad
  \nu_1, \nu_2 \geqslant 1, \quad \nu_1 = \nu_2 = 1
\]
$d_1 = d_2 = N$, \emph{$N$ quelconque}.

4) $\operatorname{Card} I = 3$ : \struck{$\frac{2}{N} = \ldots$}
\[
  \frac{1}{d_1} + \frac{1}{d_2} + \frac{1}{d_3} > 1 \qquad
  \Bigl( \frac{2}{N} = \frac{1}{p} + \frac{1}{q} + \frac{1}{r} - 1 \Bigr)
\]
\textbf{Obs.} \struck{$d_1 \geqslant d_2 \geqslant d_3$}

a) $r = 2$ :
$\frac{1}{p} + \frac{1}{q} > \frac{1}{2}$ i.e., posant $p = \alpha + 2$,
$q = \beta + 2$ $(\alpha \geqslant \beta \geqslant 0)$,
\[
  2p + 2q - pq > 0
\]
\struck{\ill{}}
\[
  \underbrace{2(\alpha + 2) + 2(\beta + 2) - (\alpha + 2)(\beta + 2)}_{4 - \alpha\beta} > 0
\]
$\alpha\beta < 4$ :
$\beta = 0$, $\alpha$ quelc. ; $\beta = 1$, $\alpha = 1, 2, 3$ ;
\struck{$\beta = 2$ \ill{}}
\[
  \begin{cases}
    (2, 2, p) & p \geqslant 2 \text{ quelc.} \qquad N = 2p \\
    (2, 3, 3) & N = 12 \\
    (2, 3, 4) & N = 24 \\
    (2, 3, 5) & N = 60
  \end{cases}
\]

b) $r \geqslant 3$ :
$\frac{1}{p} + \frac{1}{q} > 1 - \frac{1}{r} \geqslant \frac{2}{3}$
\emph{impossible} car $\frac{1}{p}, \frac{1}{q} \leqslant \frac{1}{3}$.

5) $\operatorname{Card} I \geqslant 4$. \marginal{Faisons d'abord $\operatorname{Card} I = 4$}
\[
  \frac{1}{p} + \frac{1}{q} + \frac{1}{r} + \frac{1}{s} > 2
  \qquad p \geqslant q \geqslant r \geqslant s
\]
\note{devant $p \geqslant q$ un signe biffé ; après $r$ une surcharge.}
\struck{$\frac{1}{p} + \frac{1}{q} + \frac{1}{r} > 2 - \frac{1}{s} = 1 + \bigl(1 - \frac{1}{s}\bigr)$}
\struck{$\geqslant 1 + \bigl(1 - \frac{1}{p}\bigr) > 1$}
\note{ligne biffée d'un trait ; une première formule, biffée, précède aussi $\frac{1}{p}$ dans la ligne du dessus.}

\page{67}
\note{toute la page est barrée de grands traits obliques ; elle se lit néanmoins en grande partie. Elle poursuit le cas $\operatorname{Card} I \geqslant 4$ de la p.~65.}
$(p, q, r)$ dans cet ordre, parmi les \struck{\ill{}} \uncertain{syst. admissibles} :

a) $p, 2, 2$ :
$\frac{1}{p} + \frac{1}{2} + \frac{1}{2} > 1 + \bigl(1 - \frac{1}{s}\bigr)$
i.e. $\frac{1}{p} + \frac{1}{s} > 1$ \emph{impossible} car
$\frac{1}{p}, \frac{1}{s} \leqslant \frac{1}{2}$.

b) $(p, 3, 2)$ $(p = 3, 4, 5$ ; $s \geqslant p)$
\struck{$\frac{1}{p} + \frac{1}{s} + \frac{1}{2} \ldots$}
\[
  \frac{1}{s} + \frac{1}{p} + \underbrace{\frac{1}{3} + \frac{1}{2}}_{5/6} > 2
  \qquad \frac{1}{s} + \frac{1}{p} > \frac{7}{6}
\]
impossible car $\frac{1}{s}, \frac{1}{p} \leqslant \frac{1}{3}$.
\[
  \frac{1}{p} + \frac{1}{q} > 2 - \underbrace{\Bigl(\frac{1}{r} + \frac{1}{s}\Bigr)}_{\leqslant 1} \geqslant 1
  \quad\text{donc}\quad \frac{1}{p} + \frac{1}{q} \geqslant 1
\]
\note{le premier $p$ se lit aussi bien $r$ ; sous l'accolade : « car $\frac{1}{r} \leqslant \frac{1}{2}$, $\frac{1}{s} \leqslant \frac{1}{2}$ ».}
impossible car $\frac{1}{p}, \frac{1}{q} \leqslant \frac{1}{2}$.

\emph{Impossible} !

Cas où $\operatorname{Card} I \geqslant 4$, \ldots
\[
  \frac{1}{d_1} + \frac{1}{d_2} + \frac{1}{d_3} + \frac{1}{d_4}
  > 2 + \sum_{i \geqslant 5} \Bigl(1 - \frac{1}{d_i}\Bigr)
\]
\note{après $\frac{1}{d_4}$, un terme biffé ; la borne supérieure de la somme est illisible.}
et on voit de même que c'est impossible !
\marginal{v.~2}

\emph{Étude du cas} \struck{\ill{}} \add{$(N, N)$} :
En caractéristique $0$, ou si \ill{}, on trouve
\[
  G \simeq \mu_N(k) \simeq \mathbb{Z}/N\mathbb{Z} \ldots
\]
\struck{Plus} \ill{}, il \ill{} $p$ premier à $N$
(cas de \uncertain{car.}\ $p$, \ill{} \uncertain{séparable}),
si $p \mid N$, on trouve une \add{$p$-}extension (d'un $N'$-cyclique)
\ill{} $(N = p^{r} N'$, $(p, N') = 1)$. C'est donc une
\uncertain{extension} \struck{est} \uncertain{résoluble},
\add{mais mieux :} \struck{\ill{}} \add{les}
\struck{$p$-extensions} \ill{}
\struck{\ill{} $H/$\ill{} \ill{}} \add{\ill{}}
\uncertain{sous-groupes} de $GP(1, k)$ \struck{\ill{}} abéliens et contenus dans
\ill{} conjugués entre eux, et isomorphes à $(\mathbb{Z}/p\mathbb{Z})^{r}$.
\struck{abéliens unipotents} sous-groupes unipotents de dim $1$, on trouve donc
\[
  G \subset \text{Borel},
\]
et par suite $G$ \uncertain{produit semi-direct} d'un groupe cyclique premier
à $p$ \ill{}, par un $(\mathbb{Z}/p\mathbb{Z})^{r}$. En fait
\ill{} \struck{\ill{}} (\ill{} dans $G$ premier à $p$ \ill{}
premier : \ill{}), on \ill{} \ill{}.
\marginal{En tout cas \ill{} extension résoluble
\uncertain{(d'ordre premier à $p$)} par un $(\mathbb{Z}/p^{r}\mathbb{Z})$}
\note{marge gauche, cinq lignes accolées au passage précédent ; l'exposant final est douteux.}

\emph{Étude du cas} $(N, 2, 2)$, \uncertain{$|G| = 2N$}.
En car.\ $\neq 2$, l'extension \add{\ill{}} \uncertain{de ramification} $2$ est
\uncertain{à ramification modérée}, et faisons une extension kummérienne de
degré $2$, on trouve un $X'/Y'$ qui n'a plus que $2$ pts de ramification
\struck{\ill{}}, donc $X'/Y'$ est justiciable de l'étude précédente
\ill{} \struck{\ill{}} (\ill{} $X'$ de genre $0$ \ill{}).

\begin{tikzcd}
  X \arrow[d] & X' \arrow[l] \arrow[d] \\
  Y & Y' \arrow[l]
\end{tikzcd}

\section{Paradigme \uncertain{« gruppentheoretisch »} des surfaces compactes à opérateurs}
\note{titre de sa main sur un feuillet par ailleurs blanc (p.~69) ; « gruppentheoretisch » est écrit au-dessus d'un mot biffé, peut-être « topologique ».}

\page{70}
\textbf{Paradigme des courbes algébriques}

\textbf{I} \quad \emph{Point de vue topologique.}

On considère des surfaces topologiques (sans bord) orientables
paracompactes $X$ -- on se donne aussi un ensemble $\omega$ de cardinal $2$,
et on pose
\[
  T = \mathbb{Z} \wedge_{\mathbb{Z}/2} \omega ,
\]
et on suppose donnée pour $X$ une $\omega$-orientation, i.e. une
\add{$\mathbb{Z}/2$-}application \struck{\ill{}} \add{dans}
l'ens.\ des orientations de $X$, ou ce qui revient au même, un isomorphisme
du faisceau des entiers tordus sur $X$ avec le faisceau constant $T_X$.
On suppose que les composantes connexes $X_i$ de $X$ \struck{sont} \add{est}
homéomorphes à un $\hat{X}_i \smallsetminus S_i$, où $\hat{X}_i$ est une
surface compacte, et $S_i$ une partie finie de $\hat{X}_i$. On sait que
$(\hat{X}_i, S_i)$ est déterminé à isom.\ unique près par $X_i$
(comme le décompactifié de Čech

\page{71}
de $X_i$ \uncertain{rajoutant} \uncertain{son point à l'infini}\ldots).
\uncertain{Ainsi} $\hat{X}_i$ est automatiquement $\omega$-orienté, donc aussi
\[
  \hat{X} = \coprod \hat{X}_i .
\]
De cette façon, la catégorie envisagée (des surfaces topologiques
« spéciales » \add{\struck{\ill{}} munies d'une $\omega$-orientation}
\struck{$\omega$-orientation}) devient équivalente à \add{$\omega$-orientées}
celle des surfaces topologiques à comp.\ conn.\ compactes \emph{munies}
d'une partie \struck{\ill{}} $S$ discrète, \ill{} les \uncertain{isom.}\ de
couples $(\hat{X}, S)$.

De même, la catégorie des couples \struck{$(X,$ \ill{}$)$} $(X, o, J)$ où
$(X, o)$ comme \uncertain{dessus} ($o$ la $\omega$-orientation) et $J$ partie
discrète de $X$, est équivalente à celle des systèmes
\[
  (\hat{X}, o, S, J),
\]
où $(\hat{X}, o)$ est une surface $\omega$-orientée à comp.\ conn.\ compactes,

\page{72}
et $S$, $J$ sont deux parties discrètes disjointes de $\hat{X}$.

On va passer aux catégories isotopiques correspondantes, et en donner des
interprétations en termes de théorie des groupes. Bien entendu, pour cette
description, on est ramené au cas connexe. On peut \struck{aussi} noter
\uncertain{tt de suite} que pour ce pb, on est ramené au cas où
$J = \varnothing$ [en \struck{remplaçant} \add{donnant} dans le cas général
comme autre \uncertain{donnée} celle d'un $(\hat{X}, I)$ \emph{plus} la donnée
d'une partie $J$ de $I$]. Mais on gardera le pb sous sa forme générale, en
gardant également la dissymétrie dans le rôle de $S$, $J$ -- mais

\page{73}
on devra expliciter dans le paradigme le passage d'une description à
l'autre.

\struck{La description} \add{\uncertain{$S \neq \varnothing$},}

1) Cas $J = \varnothing$. On va supposer, pour notre description, qu'on
n'est pas dans le cas $g = 0$, $\nu = 0, 1$
($g$ = genre de $\hat{X}$, $\nu = \operatorname{card} S$).
On associe à $(\hat{X}, \hat{o}, S)$ \add{$(\simeq (X, o))$} le groupe
fondamental $\Pi = \pi_1(X)$, considéré comme groupe extérieur, mais muni de
sa $T$-structure : \uncertain{les} lacets indexés par $S$, i.e. une
\struck{famille} \add{application}
\[
  S \longrightarrow \operatorname{Homext}(T, \Pi) \quad (4^{\circ})
\]
\note{« 4° » est cerclé au-dessus de la flèche ; une première écriture de la flèche, biffée, est dessous.}
Ceci permet \struck{\ill{}} $(s \in S)$ de définir
$\Pi_{\wedge}$ $(\simeq \pi_1(\hat{X}))$ comme quotient de $\Pi$ par le
sous-groupe invariant engendré par \uncertain{les images des}
$T \to \Pi$ \ill{} \ldots\ \uncertain{si $\nu = 0$ i.e. $S = \varnothing$}
\ill{}.
Ceci donne une \ill{}
\struck{\ill{}} \uncertain{-isotopique}, et \ill{} \uncertain{-isotopique},
\uncertain{sauf} si $g = 0$, $\nu = 2$ et $g = 1$, $\nu = 0$, où
la composante \struck{neutre} \ill{}
\marginal{\uncertain{plus haut} \ill{} $T \to [\ (\Pi_{\mathrm{ab}})^{2}\ ]$
\ill{} \uncertain{T identifiant} \ill{} $\Pi_{\wedge}$ \ldots}
\note{la marge gauche est écrite en travers, en plusieurs couches ; on n'en donne que les fragments sûrs.}

\page{74}
du groupe des autom.\ de $(X, o)$ n'est pas contractile (mais un
$K(\mathbb{Z}, 1)$ resp.\ un $K(\mathbb{Z}^{2}, 1)$\ldots).

2°) $J \neq \varnothing$. On a une description plus jolie (\uncertain{tous}
exception faite du cas $g = 0$, $\nu = 0, 1$) par le groupoïde fondamental de
$X$ en $J$, mais avec $T$-structure : lacets sur le groupe extérieur $\Pi$
associé -- et si $\nu = 0$, à la place de cette structure, il faut se donner
\[
  T \hookrightarrow \Bigl( \bigwedge\nolimits^{2} \Pi_{\wedge\,\mathrm{ab}} \Bigr)^{\vee}
\]
comme dans 1°\ldots. \uncertain{Ceci} \uncertain{on trouve}
une description $\infty$-isotopique, même dans le cas $g = 0$, $\nu = 2$,
ou $g = 1$, $\nu = 0$.
\marginal{\uncertain{exception}}

\section{\uncertain{Quelques problèmes spéciaux}}
\note{mots de sa main, en haut à droite d'un feuillet par ailleurs blanc (p.~75).}

\page{76}
\drawing{deux diagrammes en éventail. À gauche : $R$ au sommet, quatre
flèches $p\colon R \to S$, $q'\colon R \to A$, $q''\colon R \to A$,
$r\colon R \to F$. À droite, le même avec $R^{\circ} = R$,
$p^{\circ} = p$, $q'^{\circ}$, $q''^{\circ}$ et
$q^{\circ} = q \circ \sigma_2$ ; sous $A$, « $q' \circ \sigma_2$ » ; les
indications sur les flèches du milieu sont surchargées.}

\drawing{en dessous à gauche, un croquis de triangles accolés portant les
sommets $s$, $s'$, $s''$, $s_1$, $\bar{s}$ et les faces $f$, $f'$, $f''$,
$f_1$, $\bar{f}$, $\bar{f}_1$, $a$, $b$ ; une flèche courbe renvoie au texte.}
\[
  E = (R, q') \times_A (R, q'')
  = \bigl\{ \bigl((s, f), (\bar{s}, \bar{f})\bigr) \bigm|
    q'(s, f) = q''(\bar{s}, \bar{f}) \bigr\}
\]
\note{au-dessus des deux couples, une accolade : « repère $(s, a, \bar{f})$ ».}

\begin{tikzcd}
  & E \arrow[dl, "\pi_{01}"'] \arrow[d, "\pi_{02}"] \arrow[dr, "\pi_{12}"] & \\
  R & R & R
\end{tikzcd}

\[
  \begin{aligned}
    \pi_{01}^{+}(s, f, \bar{s}, \bar{f}) &= (s, f) \\
    \pi_{12}^{+}(s, f, \bar{s}, \bar{f}) &= (\bar{s}, \bar{f}) \\
    \pi_{02}^{+}(s, f, \bar{s}, \bar{f}) &= (s, \bar{f})
  \end{aligned}
\]
\[
  \begin{aligned}
    p(s, f) &= s &\qquad \sigma_0(s, f) &= (s', f') \\
    r(s, f) &= f & \sigma_2(s, f) &= (s'', f'') \\
    q'(s, f) &= a & & \\
    q''(s, f) &= b & \sigma_0 \sigma_2(s, f) &= \sigma_2 \sigma_0(s, f) = (s'', f'')
  \end{aligned}
\]
\note{dans la dernière égalité, les primes des arguments sont incertains.}
\[
  E^{+} = \bigl\{ (s, f, \bar{s}, \bar{f}) \in E \bigm| (s, a, \bar{f})
  \ \text{direct} \bigr\}
\]
\[
  a = q'(s, f) = q''(\bar{s}, \bar{f})
\]
\[
  \begin{cases}
    \sigma_0(s, f, \bar{s}, \bar{f}) = (s, f_1, \bar{s}, \bar{f}) \\
    \sigma_1(s, f, \bar{s}, \bar{f}) = (s, f_1, \bar{s}', \bar{f}) \\
    \sigma_2(s, f, \bar{s}, \bar{f}) = (s, f, \bar{s}_1, f_1)
  \end{cases}
\]
\note{ces trois lignes sont surchargées ; les indices $f_1$, $\bar{s}'$, $\bar{s}_1$ sont des lectures douteuses.}

\begin{tikzcd}
  & E^{+} \arrow[dl, "\pi_{01}^{+}"'] \arrow[d, "\pi_{02}^{+}"] \arrow[dr, "\pi_{12}^{+}"] & \\
  R & R & R \\
  & E^{-} \arrow[ul, "\pi_{01}^{-}"] \arrow[u, "\pi_{02}^{-}"'] \arrow[ur, "\pi_{12}^{-}"'] &
\end{tikzcd}

\note{à droite de ce diagramme, $E^{+}$ et $E^{-}$ sont écrits une seconde fois, sans flèches.}

\page{77}
\note{page de figures au crayon, sans texte suivi.}
\drawing{la double pyramide de la p.~63, redessinée :
$E^{+}$ en haut (avec, à côté, « $\rho_0 =$ » laissé en suspens), $E^{-}$ en
bas, l'équateur $D_{20}$, $D_{01}$, $D_{12}$, les points $S_0$, $S_1$, $S_2$
et les flèches $\pi_0^{\pm}$, $\pi_1^{\pm}$, $\pi_2^{\pm}$.}
\drawing{à droite, un pavage par des pentagones ; dans l'un d'eux, un petit
triangle subdivisé aux sommets $s_0$, $s_1$, $s_2$, milieux $d_{01}$,
$d_{12}$, $d_{20}$ et centre $e^{+}$ ; une flèche courbe y renvoie.}
\drawing{en bas, une sphère avec un parallèle, puis un cercle coupant une
droite graduée $-1$, $0$, $\frac{1}{2}$, $1$, $2$, les arcs passant par un
point marqué $j$ au-dessus et $\bar{j}$ au-dessous de l'axe (signes
douteux) ; à gauche, $\mathbb{P}^{1}_{\mathbb{Q}} \xleftarrow{\ \varphi\ } \mathbb{P}^{1}_{\mathbb{Q}}$.}

\page{79}
\textbf{Groupe profini \emph{spécial}}

C'est un système
$\Sigma = \bigl(T, \Pi, I, (\Phi_i)_{i \in I} \in (\operatorname{Homext}(T, \Pi))^{I}\bigr)$
avec
\[
  \Sigma
  \begin{cases}
    T \text{ groupe profini} \simeq \hat{\mathbb{Z}} \\
    \Pi \text{ groupe profini} \\
    I \text{ ensemble fini} \\
    \Phi_i \in \operatorname{Homext}(T, \Pi) \quad \text{pour } i \in I
  \end{cases}
\]
avec l'axiome que ce système est \struck{\ill{}} \add{isomorphe} à l'un des
systèmes $\Sigma_\nu$ $(\nu \in \mathbb{N}, \nu \geqslant 1)$ suivants :
\[
  \Sigma_\nu
  \begin{cases}
    T = \hat{\mathbb{Z}} \\
    \Pi = \text{groupe profini } \Pi_\nu = \widehat{L_\nu}, \text{ où }
      L_\nu = \{\lambda_0, \ldots, \lambda_\nu \mid \lambda_0 \lambda_1 \cdots \lambda_\nu = 1\} \\
    \qquad \text{(donc $\Pi$ profini libre : $\nu$ générateurs)} \\
    I = [0, \nu]_{\mathbb{N}} \\
    \Phi_i = \text{classe de l'hom. } \hat{\mathbb{Z}} \xrightarrow{\varphi_i} \Pi
      \text{ tel que } \varphi_i(1) = \lambda_i
  \end{cases}
\]
\note{la relation de $L_\nu$ est coupée par le bord de la feuille ; « $= 1$ » est suppléé par le sens.}

\textbf{N.B.} $I$ peut être interprété comme une partie de
$\operatorname{Homext}(T, \Pi)$. Si $\Sigma \simeq \Sigma_\nu$, $\nu$ est
uniquement déterminé comme $\operatorname{card} I - 1$, ou
\struck{\ill{}} \uncertain{aussi} comme rang profini de $\Pi_{\mathrm{ab}}$.
Pour $\nu$ fixé, la catégorie des $\Sigma$ correspondants est équivalente à la
catégorie des torseurs sous
\[
  \widetilde{A}_\nu = \operatorname{Aut}(\Sigma_\nu) .
\]
On a des homomorphismes
\[
  \widetilde{A}_\nu \longrightarrow \mathfrak{S}_{\nu+1} \quad (= \operatorname{Aut}(I_\nu)),
  \qquad
  \widetilde{A}_\nu \xrightarrow{\ \chi\ } \operatorname{Aut} \hat{\mathbb{Z}} \simeq \hat{\mathbb{Z}}^{*}
\]
\note{sous $\mathfrak{S}$, un indice surchargé ; on lit $\nu + 1$ comme à la page suivante.}

\page{80}
donnant naissance :
\[
  \widetilde{A}_\nu \longrightarrow \mathfrak{S}_{\nu+1} \times \hat{\mathbb{Z}}^{*}
\]
qui est \emph{surjectif}. Le noyau de $\widetilde{A}_\nu \to \mathfrak{S}_{\nu+1}$
se note $A_\nu$, celui de
$\widetilde{A}_\nu \to \mathfrak{S}_{\nu+1} \times \hat{\mathbb{Z}}^{*}$
se note $A_\nu^{\circ}$. \uncertain{J'affirme} un
homomorphisme injectif
\[
  \Pi \longrightarrow A_\nu^{\circ} \subset A_\nu
\]
via automorphismes intérieurs. Le conoyau de $\Pi \to A_\nu^{\circ}$
($\Pi \to A_\nu$, \struck{$\Pi \to \widetilde{A}_\nu$})
$\bigl(\Gamma_\nu, \widetilde{\Gamma}_\nu\bigr)$ se note $\Gamma_\nu^{\circ}$,
\uncertain{d'où} \ill{} des suites exactes
\[
  \begin{cases}
    1 \to \Pi \to A_\nu^{\circ} \to \Gamma_\nu \to 1 \\
    1 \to A_\nu \to \widetilde{A}_\nu \to \mathfrak{S}_{\nu+1} \to 1 \\
    1 \to A_\nu^{\circ} \to A_\nu \to \hat{\mathbb{Z}}^{*} \to 1
  \end{cases}
\]
\note{entre la première et la deuxième ligne, une ligne biffée, illisible ; dans la première, $\Gamma_\nu$ sans l'exposant $\circ$ annoncé plus haut.}
\[
  \overbrace{\Pi \quad \Gamma_\nu^{\circ} \quad \hat{\mathbb{Z}}^{*} \quad \mathfrak{S}_{\nu+1}}^{\widetilde{A}_\nu}
  \qquad \text{filtration}
\]

\emph{Exemples} (1) Soit $X$ droite projective sur
$k$ alg.\ clos de car.\ $0$, $I \subset X(k)$, $\operatorname{card} I = \nu + 1$,
$T = T(k) = \varprojlim_n \mu_n(k)$, $\Pi = \pi_1(X \smallsetminus I, x)$
où $x \in X(k) - I$. Il y a \ill{} \ldots
\note{l'exemple se poursuit au-delà de ce lot.}

\end{document}
